\documentclass[10pt,aspectratio=169]{beamer}

% ------------------------------------------------------------------
% Presentació adaptada: Descriptiva univariant amb orientació computacional
% Grau en Enginyeria Informàtica - Estadística
% ------------------------------------------------------------------

\usepackage[utf8]{inputenc}
\usepackage[T1]{fontenc}
\usepackage[catalan]{babel}
\usepackage{amsmath,amssymb}
\usepackage{booktabs}
\usepackage{array}
\usepackage{tikz}
\usepackage{pgfplots}
\usepgfplotslibrary{statistics}
\usepackage{xcolor}
\usepackage{listings}
\usepackage{eurosym}
\pgfplotsset{compat=1.18}

% ------------------------------------------------------------------
% Estil visual
% ------------------------------------------------------------------
\definecolor{UABGreen}{RGB}{0,103,71}
\definecolor{UABLight}{RGB}{224,241,236}
\definecolor{DeepBlue}{RGB}{20,65,110}
\definecolor{AccentOrange}{RGB}{230,126,34}
\definecolor{SoftGray}{RGB}{245,247,250}
\definecolor{CodeBack}{RGB}{248,248,248}
\definecolor{CodeFrame}{RGB}{210,215,220}

\usetheme{Madrid}
\usecolortheme[named=UABGreen]{structure}
\setbeamercolor{title}{fg=white,bg=UABGreen}
\setbeamercolor{frametitle}{fg=white,bg=UABGreen}
\setbeamercolor{block title}{fg=white,bg=DeepBlue}
\setbeamercolor{block body}{fg=black,bg=SoftGray}
\setbeamercolor{alerted text}{fg=AccentOrange}
\setbeamertemplate{navigation symbols}{}
\setbeamertemplate{itemize item}{\color{UABGreen}$\blacktriangleright$}
\setbeamertemplate{itemize subitem}{\color{AccentOrange}$\bullet$}
\setbeamertemplate{enumerate item}{\color{UABGreen}\insertenumlabel.}

\setbeamertemplate{footline}{%
  \leavevmode%
  \hbox{%
  \begin{beamercolorbox}[wd=.72\paperwidth,ht=2.4ex,dp=1ex,leftskip=1.4em]{author in head/foot}%
    Estadística -- Grau en Enginyeria Informàtica
  \end{beamercolorbox}%
  \begin{beamercolorbox}[wd=.28\paperwidth,ht=2.4ex,dp=1ex,center]{date in head/foot}%
    \insertframenumber{} / \inserttotalframenumber
  \end{beamercolorbox}}%
  \vskip0pt%
}

\lstdefinelanguage{R}{%
  morekeywords={if,else,repeat,while,function,for,in,next,break,TRUE,FALSE,NULL,NA,NaN,Inf,
    library,mean,median,sd,var,summary,table,prop.table,quantile,ggplot,aes,geom_histogram,
    geom_boxplot,summarise,group_by,mutate,read_csv,filter,arrange,set.seed,runif,rnorm,
    future_map_dbl,plan,multisession},
  sensitive=true,
  morecomment=[l]\#,
  morestring=[b]",
  morestring=[b]'
}
\lstset{%
  language=R,
  basicstyle=\ttfamily\scriptsize,
  keywordstyle=\color{DeepBlue}\bfseries,
  commentstyle=\color{gray},
  stringstyle=\color{AccentOrange},
  backgroundcolor=\color{CodeBack},
  frame=single,
  rulecolor=\color{CodeFrame},
  breaklines=true,
  columns=fullflexible,
  keepspaces=true,
  showstringspaces=false,
  literate={à}{{\`a}}1 {è}{{\`e}}1 {é}{{\'e}}1 {í}{{\'i}}1 {ï}{{\"i}}1 {ò}{{\`o}}1 {ó}{{\'o}}1 {ú}{{\'u}}1 {ü}{{\"u}}1 {ç}{{\c{c}}}1 {À}{{\`A}}1 {È}{{\`E}}1 {É}{{\'E}}1 {Í}{{\'I}}1 {Ò}{{\`O}}1 {Ó}{{\'O}}1 {Ú}{{\'U}}1
}

\newcommand{\R}{\textsf{R}}
\newcommand{\RStudio}{\textsf{RStudio/Posit}}
\newcommand{\GitHub}{\textsf{GitHub}}
\newcommand{\important}[1]{\textcolor{UABGreen}{\textbf{#1}}}
\newcommand{\exemple}[1]{\textcolor{AccentOrange}{\textbf{#1}}}

% ------------------------------------------------------------------
% Dades de la presentació
% ------------------------------------------------------------------
\title[Descriptiva univariant computacional]{Estadística descriptiva univariant}
\author{Estadística -- Grau en Enginyeria Informàtica}
\date{Curs 2026--27}

\begin{document}

% ------------------------------------------------------------------
\begin{frame}[plain]
  \titlepage
  \vspace{-0.5cm}
  \begin{center}
  \small
  \textcolor{DeepBlue}{Dades, codi, visualització, simulació i reproduïbilitat}
  \end{center}
\end{frame}

% ------------------------------------------------------------------
\begin{frame}{Per què començar per descriptiva univariant?}
\begin{columns}[T]
\column{0.55\textwidth}
\begin{itemize}
  \item Abans de modelitzar, predir o programar un algoritme, cal \important{mirar les dades}.
  \item La descriptiva univariant respon preguntes bàsiques:
  \begin{itemize}
    \item Quins valors apareixen?
    \item Quin és un valor típic?
    \item Hi ha molta variabilitat?
    \item Hi ha valors estranys o errors?
  \end{itemize}
  \item En informàtica això apareix en \exemple{logs}, \exemple{benchmarks}, \exemple{latència}, \exemple{errors}, \exemple{memòria} i \exemple{experiments A/B}.
\end{itemize}

\column{0.4\textwidth}
\begin{block}{Idea clau}
Una variable no és només una columna: és una història sobre un sistema.
\end{block}
\vspace{0.2cm}
\begin{tikzpicture}
\node[draw=UABGreen,rounded corners,fill=UABLight,text width=4.3cm,align=center,inner sep=8pt] {
Dades $\rightarrow$ Resum $\rightarrow$ Gràfic $\rightarrow$ Interpretació $\rightarrow$ Decisió
};
\end{tikzpicture}
\end{columns}
\end{frame}

% ------------------------------------------------------------------
\begin{frame}{Objectius del tema}
Al final d'aquest tema hauríeu de poder:
\begin{enumerate}
  \item Reconèixer el tipus d'una variable i escollir el resum adequat.
  \item Construir taules de freqüències i gràfics amb \R.
  \item Calcular i interpretar mitjana, mediana, moda, variància, desviació típica, quartils i percentils.
  \item Detectar asimetria, valors atípics i possibles errors de registre.
  \item Preparar un petit informe reproduïble amb \RStudio, Quarto i \GitHub.
\end{enumerate}
\vspace{0.2cm}
\begin{block}{Enfocament computacional}
No farem només càlculs manuals: aprendrem a \important{automatitzar}, \important{visualitzar} i \important{documentar} l'anàlisi.
\end{block}
\end{frame}

% ------------------------------------------------------------------
\begin{frame}{El flux de treball que farem servir}
\begin{center}
\begin{tikzpicture}[node distance=1.55cm, every node/.style={font=\small}]
\node[draw,rounded corners,fill=UABLight,minimum width=2.3cm,minimum height=0.8cm] (data) {Dades};
\node[draw,rounded corners,fill=UABLight,right of=data,xshift=1.6cm,minimum width=2.3cm,minimum height=0.8cm] (code) {Codi R};
\node[draw,rounded corners,fill=UABLight,right of=code,xshift=1.6cm,minimum width=2.3cm,minimum height=0.8cm] (plot) {Gràfics};
\node[draw,rounded corners,fill=UABLight,right of=plot,xshift=1.6cm,minimum width=2.3cm,minimum height=0.8cm] (report) {Informe};
\draw[->,thick,UABGreen] (data) -- (code);
\draw[->,thick,UABGreen] (code) -- (plot);
\draw[->,thick,UABGreen] (plot) -- (report);
\end{tikzpicture}
\end{center}

\vspace{0.4cm}
\begin{columns}[T]
\column{0.48\textwidth}
\begin{block}{Eines}
\begin{itemize}
  \item \R{} per calcular.
  \item \RStudio{} per treballar per projectes.
  \item Quarto per fer informes.
  \item \GitHub{} per versionar i lliurar.
\end{itemize}
\end{block}
\column{0.48\textwidth}
\begin{block}{Bones pràctiques}
\begin{itemize}
  \item Codi curt i llegible.
  \item Dades separades del codi.
  \item Resultats regenerables.
  \item Commits amb missatges informatius.
\end{itemize}
\end{block}
\end{columns}
\end{frame}

% ------------------------------------------------------------------
\section{Variables i dades}

\begin{frame}{Tipus de variables: exemples informàtics}
\begin{columns}[T]
\column{0.48\textwidth}
\begin{block}{Qualitatives}
\begin{itemize}
  \item \important{Nominals}: sistema operatiu, navegador, estat d'una petició \texttt{OK/ERROR}, tipus d'usuari.
  \item \important{Ordinals}: prioritat \texttt{baixa/mitjana/alta}, severitat d'un bug, nivell de risc.
\end{itemize}
\end{block}
\column{0.48\textwidth}
\begin{block}{Quantitatives}
\begin{itemize}
  \item \important{Discretes}: nombre d'errors, peticions per minut, commits per setmana.
  \item \important{Contínues}: latència, temps d'execució, ús de memòria, temperatura d'un sensor.
\end{itemize}
\end{block}
\end{columns}
\vspace{0.35cm}
\begin{alertblock}{Pregunta recurrent}
El fet que una variable estigui codificada amb números no vol dir necessàriament que sigui quantitativa.
\end{alertblock}
\end{frame}

% ------------------------------------------------------------------
\begin{frame}{La matriu de dades}
\begin{columns}[T]
\column{0.48\textwidth}
\begin{itemize}
  \item Les \important{files} són casos o observacions.
  \item Les \important{columnes} són variables.
  \item La mida mostral és el nombre de files: $n$.
  \item Cada columna ha de tenir un significat i unes unitats clares.
\end{itemize}

\column{0.48\textwidth}
\centering
\small
\begin{tabular}{llllr}
\toprule
id & navegador & estat & mida & temps\_ms \\
\midrule
1 & Chrome  & OK    & gran  & 118 \\
2 & Firefox & OK    & petit & 74 \\
3 & Chrome  & ERROR & mitjà & 312 \\
4 & Safari  & OK    & gran  & 141 \\
5 & Chrome  & OK    & petit & 69 \\
\bottomrule
\end{tabular}
\vspace{0.2cm}

\scriptsize Exemple: registre simplificat de peticions a una aplicació web.
\end{columns}
\end{frame}

% ------------------------------------------------------------------
\begin{frame}[fragile]{Llegir dades amb R}
\begin{lstlisting}
library(tidyverse)

logs <- read_csv("data/logs_api.csv")

glimpse(logs)
summary(logs$temps_ms)
\end{lstlisting}
\vspace{0.3cm}
\begin{block}{Objectiu}
Passar d'un fitxer de dades a un objecte que podem explorar, resumir i representar gràficament.
\end{block}
\end{frame}

% ------------------------------------------------------------------
\section{Freqüències i gràfics}

\begin{frame}{Freqüència absoluta i relativa}
Per una variable amb valors $x_1,\ldots,x_k$:
\vspace{0.25cm}
\begin{columns}[T]
\column{0.45\textwidth}
\begin{block}{Freqüència absoluta}
\[
 n_i = \text{nombre d'observacions amb valor } x_i
\]
\[
 \sum_{i=1}^k n_i = n
\]
\end{block}
\column{0.45\textwidth}
\begin{block}{Freqüència relativa}
\[
 f_i = \frac{n_i}{n}
\]
\[
 \sum_{i=1}^k f_i = 1
\]
\end{block}
\end{columns}
\vspace{0.25cm}
\begin{alertblock}{Interpretació}
Les freqüències relatives permeten comparar mostres de mides diferents.
\end{alertblock}
\end{frame}

% ------------------------------------------------------------------
\begin{frame}[fragile]{Taules de freqüències amb R}
\begin{columns}[T]
\column{0.48\textwidth}
\begin{lstlisting}
# Freqüències absolutes
freq_abs <- table(logs$estat)

# Freqüències relatives
freq_rel <- prop.table(freq_abs)

freq_abs
freq_rel
\end{lstlisting}
\column{0.48\textwidth}
\small
\begin{tabular}{lrr}
\toprule
estat & $n_i$ & $f_i$ \\
\midrule
OK    & 824 & 0.824 \\
ERROR & 176 & 0.176 \\
\bottomrule
\end{tabular}
\vspace{0.25cm}
\begin{block}{Pregunta}
És acceptable una taxa d'error del 17.6\%? Depèn del context i del sistema.
\end{block}
\end{columns}
\end{frame}

% ------------------------------------------------------------------
\begin{frame}{Representació de variables qualitatives}
\begin{columns}[T]
\column{0.5\textwidth}
\begin{itemize}
  \item Per variables qualitatives fem servir sobretot \important{diagrames de barres}.
  \item El diagrama de sectors pot ser útil amb poques categories, però sovint és menys informatiu.
  \item En informàtica interessa veure ràpidament quina categoria domina.
\end{itemize}
\column{0.46\textwidth}
\begin{tikzpicture}
\begin{axis}[
    ybar,
    width=6.2cm,
    height=4.2cm,
    ymin=0,
    ylabel={Freqüència},
    symbolic x coords={Chrome,Firefox,Safari,Edge},
    xtick=data,
    x tick label style={rotate=25,anchor=east},
    bar width=13pt,
    nodes near coords,
    nodes near coords style={font=\scriptsize},
    enlarge x limits=0.18]
\addplot[fill=UABGreen!70] coordinates {(Chrome,430) (Firefox,260) (Safari,180) (Edge,130)};
\end{axis}
\end{tikzpicture}
\end{columns}
\end{frame}

% ------------------------------------------------------------------
\begin{frame}[fragile]{Gràfic de barres amb \texttt{ggplot2}}
\begin{lstlisting}
library(tidyverse)

ggplot(logs, aes(x = navegador)) +
  geom_bar() +
  labs(
    title = "Peticions per navegador",
    x = "Navegador",
    y = "Nombre de peticions"
  )
\end{lstlisting}
\vspace{0.15cm}
\begin{block}{Missatge metodològic}
El codi del gràfic ha de formar part de l'informe. Això permet reproduir exactament el resultat.
\end{block}
\end{frame}

% ------------------------------------------------------------------
\begin{frame}{Freqüències acumulades}
Per variables numèriques ordenables:
\[
N_i = n_1 + n_2 + \cdots + n_i, \qquad F_i = \frac{N_i}{n}.
\]
\vspace{0.3cm}
\begin{block}{Aplicació informàtica}
Si $X$ és el temps de resposta d'una API, $F(200)$ pot interpretar-se com:
\[
\text{proporció de peticions amb temps de resposta } \leq 200 \text{ ms}.
\]
\end{block}
\vspace{0.15cm}
\begin{alertblock}{Pregunta útil}
Quin percentatge de peticions compleix el llindar de qualitat del servei?
\end{alertblock}
\end{frame}

% ------------------------------------------------------------------
\begin{frame}{Histogrames: forma de la distribució}
\begin{columns}[T]
\column{0.46\textwidth}
\begin{itemize}
  \item Per variables quantitatives contínues, l'histograma resumeix la distribució.
  \item Cal escollir un nombre d'intervals o \textit{bins}.
  \item Massa pocs intervals amaguen informació; massa intervals afegeixen soroll.
\end{itemize}
\column{0.5\textwidth}
\begin{tikzpicture}
\begin{axis}[
  ybar interval,
  width=6.8cm,
  height=4.2cm,
  ymin=0,
  xlabel={temps de resposta (ms)},
  ylabel={freqüència},
  xtick={50,100,150,200,250,300,350}]
\addplot[fill=DeepBlue!65] coordinates {(50,8) (100,24) (150,35) (200,20) (250,9) (300,3) (350,0)};
\end{axis}
\end{tikzpicture}
\end{columns}
\end{frame}

% ------------------------------------------------------------------
\begin{frame}[fragile]{Histograma amb R}
\begin{lstlisting}
ggplot(logs, aes(x = temps_ms)) +
  geom_histogram(bins = 30) +
  labs(
    title = "Distribució dels temps de resposta",
    x = "Temps de resposta (ms)",
    y = "Nombre de peticions"
  )
\end{lstlisting}
\vspace{0.25cm}
\begin{block}{Experiment ràpid}
Canvieu \texttt{bins = 10}, \texttt{bins = 30} i \texttt{bins = 80}. Canvia la història que explica el gràfic?
\end{block}
\end{frame}

% ------------------------------------------------------------------
\begin{frame}{Compte amb els intervals de diferent amplada}
\begin{itemize}
  \item Si tots els intervals tenen la mateixa amplada, podem dibuixar alçades proporcionals a les freqüències.
  \item Si els intervals tenen amplades diferents, l'àrea del rectangle ha de ser proporcional a la freqüència.
  \item Altrament podem obtenir una representació visualment enganyosa.
\end{itemize}
\vspace{0.3cm}
\begin{block}{Regla pràctica}
En un histograma, el que importa és l'\important{àrea}, no només l'alçada.
\end{block}
\end{frame}

% ------------------------------------------------------------------
\section{Mesures de centre}

\begin{frame}{Mesures de centre}
\begin{columns}[T]
\column{0.33\textwidth}
\begin{block}{Mitjana}
\[
\bar{x}=\frac{1}{n}\sum_{i=1}^n x_i
\]
Valor d'equilibri. Sensible a valors extrems.
\end{block}
\column{0.33\textwidth}
\begin{block}{Mediana}
Valor central un cop ordenades les dades. Robusta davant d'extrems.
\end{block}
\column{0.33\textwidth}
\begin{block}{Moda}
Valor més freqüent. Molt útil en variables qualitatives o discretes.
\end{block}
\end{columns}
\vspace{0.25cm}
\begin{alertblock}{En dades computacionals}
Temps de resposta, mida de fitxers i ús de memòria poden tenir cues llargues. La mediana sovint és més informativa que la mitjana.
\end{alertblock}
\end{frame}

% ------------------------------------------------------------------
\begin{frame}{Mitjana i mediana amb valors extrems}
\begin{center}
\small
\begin{tabular}{lrrrrrrr}
\toprule
Execució & 1 & 2 & 3 & 4 & 5 & 6 & 7 \\
\midrule
Temps A (ms) & 98 & 101 & 99 & 100 & 102 & 97 & 103 \\
Temps B (ms) & 98 & 101 & 99 & 100 & 102 & 97 & 620 \\
\bottomrule
\end{tabular}
\end{center}
\vspace{0.2cm}
\begin{columns}[T]
\column{0.48\textwidth}
\begin{block}{Sistema A}
\[
\bar{x}=100, \qquad Md=100
\]
\end{block}
\column{0.48\textwidth}
\begin{block}{Sistema B}
\[
\bar{x}=174, \qquad Md=100
\]
\end{block}
\end{columns}
\vspace{0.25cm}
\begin{alertblock}{Interpretació}
Un sol valor extrem pot canviar molt la mitjana, però no necessàriament la mediana.
\end{alertblock}
\end{frame}

% ------------------------------------------------------------------
\begin{frame}[fragile]{Resums amb R}
\begin{lstlisting}
logs %>%
  summarise(
    n = n(),
    mitjana = mean(temps_ms),
    mediana = median(temps_ms),
    minim = min(temps_ms),
    maxim = max(temps_ms)
  )
\end{lstlisting}
\vspace{0.2cm}
\begin{block}{Bona pràctica}
No doneu només un número. Doneu també el context: unitats, mida mostral i possibles valors atípics.
\end{block}
\end{frame}

% ------------------------------------------------------------------
\begin{frame}{Quan la mitjana pot enganyar}
\begin{columns}[T]
\column{0.48\textwidth}
\begin{block}{Exemple clàssic: sous}
Quatre programadors cobren 1000\euro{} i el propietari 6000\euro{}.
\[
\bar{x}=\frac{4\cdot 1000+6000}{5}=2000
\]
\end{block}
\column{0.48\textwidth}
\begin{block}{Exemple informàtic: latència}
Moltes peticions responen en 100 ms, però algunes fallen i triguen 2 s.
\[
\text{mitjana} > \text{mediana}
\]
\end{block}
\end{columns}
\vspace{0.25cm}
\begin{alertblock}{Missatge}
Sempre que hi hagi asimetria o valors extrems, compareu mitjana i mediana.
\end{alertblock}
\end{frame}

% ------------------------------------------------------------------
\section{Mesures de dispersió}

\begin{frame}{Mesures de dispersió}
\begin{itemize}
  \item Les mesures de centre no són suficients: dos sistemes poden tenir la mateixa mitjana i variabilitat molt diferent.
  \item La dispersió indica com d'inestables o heterogènies són les observacions.
\end{itemize}
\vspace{0.25cm}
\begin{columns}[T]
\column{0.48\textwidth}
\begin{block}{Variància corregida}
\[
s^2=\frac{1}{n-1}\sum_{i=1}^n (x_i-\bar{x})^2
\]
\end{block}
\column{0.48\textwidth}
\begin{block}{Desviació típica}
\[
s=\sqrt{s^2}
\]
Té les mateixes unitats que les dades.
\end{block}
\end{columns}
\end{frame}

% ------------------------------------------------------------------
\begin{frame}{Mateixa mitjana, diferent estabilitat}
\begin{center}
\begin{tikzpicture}
\begin{axis}[
  width=10.5cm,
  height=4.7cm,
  ymin=60,ymax=150,
  xlabel={execució},
  ylabel={temps (ms)},
  legend pos=north west]
\addplot[mark=*,DeepBlue,thick] coordinates {(1,98) (2,101) (3,99) (4,100) (5,102) (6,97) (7,103) (8,100)};
\addplot[mark=square*,AccentOrange,thick] coordinates {(1,70) (2,130) (3,82) (4,145) (5,93) (6,121) (7,75) (8,144)};
\legend{Sistema estable,Sistema variable}
\end{axis}
\end{tikzpicture}
\end{center}
\begin{block}{Pregunta}
Quin sistema preferiríeu per a una aplicació en temps real?
\end{block}
\end{frame}

% ------------------------------------------------------------------
\begin{frame}[fragile]{Dispersió amb R}
\begin{lstlisting}
logs %>%
  summarise(
    mitjana = mean(temps_ms),
    variancia = var(temps_ms),
    desviacio = sd(temps_ms),
    rang = max(temps_ms) - min(temps_ms),
    cv = sd(temps_ms) / mean(temps_ms) * 100
  )
\end{lstlisting}
\vspace{0.2cm}
\begin{block}{Coeficient de variació}
\[
CV = \frac{s}{\bar{x}}\times 100
\]
Serveix per comparar variabilitat relativa entre variables o sistemes amb escales diferents.
\end{block}
\end{frame}

% ------------------------------------------------------------------
\section{Percentils, quartils i boxplots}

\begin{frame}{Percentils i quartils}
\begin{itemize}
  \item El percentil $p$ és un valor que deixa aproximadament un $p\%$ de les observacions per sota.
  \item La mediana és el percentil 50.
  \item Els quartils divideixen les dades en quatre parts:
  \[
  Q_1=P_{25}, \qquad Q_2=P_{50}=Md, \qquad Q_3=P_{75}.
  \]
  \item El rang interquartílic és
  \[
  RI = Q_3-Q_1.
  \]
\end{itemize}
\begin{alertblock}{En sistemes informàtics}
El percentil 95 o 99 de la latència pot ser més important que la mitjana.
\end{alertblock}
\end{frame}

% ------------------------------------------------------------------
\begin{frame}[fragile]{Quantils amb R}
\begin{lstlisting}
quantile(logs$temps_ms, probs = c(0.25, 0.50, 0.75, 0.95, 0.99))
\end{lstlisting}
\vspace{0.25cm}
\begin{center}
\small
\begin{tabular}{rrrrr}
\toprule
P25 & P50 & P75 & P95 & P99 \\
\midrule
82 & 106 & 148 & 310 & 870 \\
\bottomrule
\end{tabular}
\end{center}
\vspace{0.15cm}
\begin{block}{Lectura}
El 99\% de les peticions responen en 870 ms o menys. L'1\% restant pot ser crític.
\end{block}
\end{frame}

% ------------------------------------------------------------------
\begin{frame}{Diagrama de caixa: detectar valors atípics}
\begin{columns}[T]
\column{0.48\textwidth}
\begin{itemize}
  \item Resumeix mediana, quartils, rang interquartílic i possibles valors atípics.
  \item Barreres internes:
  \[
  Q_1-1.5RI, \qquad Q_3+1.5RI.
  \]
  \item Els punts fora d'aquest interval són candidats a \important{outliers}.
\end{itemize}
\column{0.48\textwidth}
\begin{tikzpicture}
\begin{axis}[
  width=6.3cm,
  height=4.6cm,
  boxplot/draw direction=y,
  xtick={1},
  xticklabels={temps},
  ylabel={ms},
  ymin=40,ymax=400]
\addplot+[boxplot prepared={lower whisker=55, lower quartile=82, median=106, upper quartile=148, upper whisker=245}, mark=*] coordinates {(1,310) (1,360)};
\end{axis}
\end{tikzpicture}
\end{columns}
\end{frame}

% ------------------------------------------------------------------
\begin{frame}[fragile]{Boxplot amb R}
\begin{lstlisting}
ggplot(logs, aes(y = temps_ms)) +
  geom_boxplot() +
  labs(
    title = "Diagrama de caixa dels temps de resposta",
    y = "Temps de resposta (ms)"
  )
\end{lstlisting}
\vspace{0.2cm}
\begin{block}{Ampliació natural}
Comparar per grups:
\begin{lstlisting}
ggplot(logs, aes(x = navegador, y = temps_ms)) + geom_boxplot()
\end{lstlisting}
\end{block}
\end{frame}

% ------------------------------------------------------------------
\section{Forma de la distribució}

\begin{frame}{Asimetria}
\begin{itemize}
  \item Una distribució és simètrica si la seva forma és semblant a banda i banda del centre.
  \item L'asimetria positiva és habitual en temps de resposta: moltes observacions petites i unes poques molt grans.
  \item L'asimetria negativa és menys habitual en latències, però pot aparèixer en notes, puntuacions o percentatges limitats superiorment.
\end{itemize}
\vspace{0.25cm}
\begin{columns}[T]
\column{0.5\textwidth}
\begin{block}{Coeficient d'asimetria}
\[
As=\frac{\frac{1}{n}\sum_{i=1}^n (x_i-\bar{x})^3}{s^3}
\]
\end{block}
\column{0.45\textwidth}
\begin{alertblock}{Senyal pràctic}
Si $\bar{x} > Md$, sovint hi ha cua a la dreta.
\end{alertblock}
\end{columns}
\end{frame}

% ------------------------------------------------------------------
\begin{frame}{Distribució simètrica i asimètrica}
\begin{center}
\begin{tikzpicture}
\begin{axis}[
  width=10.5cm,
  height=4.8cm,
  domain=-3:7,
  samples=200,
  axis lines=left,
  xlabel={$x$},
  ylabel={densitat},
  legend pos=north east,
  ymin=0]
\addplot[DeepBlue,thick] {exp(-x^2/2)/2.5066};
\addplot[AccentOrange,thick] {x>0 ? 0.8*exp(-x/1.5) : 0};
\legend{aprox. simètrica,cua a la dreta}
\end{axis}
\end{tikzpicture}
\end{center}
\begin{block}{Interpretació}
La forma de la distribució ajuda a decidir quines mesures resum són més informatives.
\end{block}
\end{frame}

% ------------------------------------------------------------------
\section{Simulació i reproduïbilitat}

\begin{frame}[fragile]{Simular per entendre}
\begin{lstlisting}
set.seed(2026)

# Temps de resposta artificials: molts valors petits i alguns grans
temps <- c(rnorm(950, mean = 100, sd = 20),
           rnorm(50, mean = 500, sd = 80))

mean(temps)
median(temps)
sd(temps)
quantile(temps, c(.50, .95, .99))
\end{lstlisting}
\vspace{0.15cm}
\begin{block}{Per què simular?}
La simulació permet construir exemples controlats i veure com canvien els estadístics quan canvia el sistema.
\end{block}
\end{frame}

% ------------------------------------------------------------------
\begin{frame}[fragile]{Primer contacte amb càlcul repetit}
\begin{lstlisting}
set.seed(2026)

resultats <- replicate(1000, {
  x <- rnorm(100, mean = 100, sd = 20)
  mean(x)
})

summary(resultats)
sd(resultats)
\end{lstlisting}
\vspace{0.2cm}
\begin{alertblock}{Idea per més endavant}
Quan repetim molts càlculs, podem plantejar-nos paral·lelitzar-los. La descriptiva és també una porta d'entrada a la simulació computacional.
\end{alertblock}
\end{frame}

% ------------------------------------------------------------------
\begin{frame}[fragile]{Exemple: versió paral·lelitzable}
\begin{lstlisting}
library(future)
library(furrr)

plan(multisession, workers = 4)

resultats <- future_map_dbl(1:1000, function(i) {
  x <- rnorm(100, mean = 100, sd = 20)
  mean(x)
})
\end{lstlisting}
\vspace{0.2cm}
\begin{block}{Missatge}
La idea estadística és la mateixa; el que canvia és com repartim el càlcul entre nuclis o processos.
\end{block}
\end{frame}

% ------------------------------------------------------------------
\begin{frame}{Reproduïbilitat amb GitHub}
\begin{columns}[T]
\column{0.48\textwidth}
\begin{block}{Estructura recomanada}
\small
\begin{tabular}{ll}
\toprule
Carpeta/fitxer & Contingut \\
\midrule
\texttt{data/} & dades originals \\
\texttt{R/} & scripts auxiliars \\
\texttt{figures/} & gràfics generats \\
\texttt{informe.qmd} & informe reproduïble \\
\texttt{README.md} & descripció del projecte \\
\bottomrule
\end{tabular}
\end{block}
\column{0.48\textwidth}
\begin{block}{Commits útils}
\begin{itemize}
  \item \texttt{afegeix lectura de dades}
  \item \texttt{calcula taula de freqüències}
  \item \texttt{afegeix histograma de latència}
  \item \texttt{interpreta valors atípics}
\end{itemize}
\end{block}
\end{columns}
\end{frame}

\end{document}
